\originalchapter{能量、动量与角动量}
\originalsection{坡印廷定理}
电场对电流的单位体积功率为 $\jj\cdot\ee$. 用麦克斯韦方程消去 $\jj$, 再用 $\nabla\cdot(\ee\times\bb)=\bb\cdot(\nabla\times\ee)-\ee\cdot(\nabla\times\bb)$, 得
\[
\jj\cdot\ee=-\partial_t\left(\frac\eps2E^2+\frac{B^2}{2\mu_0}\right)
-\nabla\cdot\left(\frac1{\mu_0}\ee\times\bb\right).
\]
定义真空场能密度 $u=\eps E^2/2+B^2/(2\mu_0)$, 坡印廷矢量 $\vct S=\ee\times\bb/\mu_0$, 则 $\partial_tu+\nabla\cdot\vct S=-\jj\cdot\ee$. 积分后是区域内场能的减少等于对物质做功加上穿过边界向外的能流. 因而电阻耗散可由侧面的电磁能流解释, 无须把能量只想成沿导线内部被电子搬运.

对静止、无损、无预存极化或磁化偏置、构成关系瞬时线性且参数不随时间改变的介质, 若构成张量对称 (实常数各向同性模型是充分情形), 可相应写 $u=(\ee\cdot\vct D+\bb\cdot\vct H)/2$, $\vct S=\ee\times\vct H$. 色散介质需把物质振子的能量并入, 不能只用静态公式决定储能正负.

\originalsection{麦克斯韦应力张量}
洛伦兹力密度是 $\vct f=\rho\ee+\jj\times\bb$. 将源用麦克斯韦方程消去,
\begin{align*}
f_i={}&\eps E_i\partial_jE_j
+\frac1{\mu_0}[(\bb\cdot\nabla)B_i-\tfrac12\partial_iB^2]
-\eps(\partial_t\ee\times\bb)_i.
\end{align*}
这里用了 $\nabla\cdot\bb=0$. 为形成场动量的时间导数, 加上 $\partial_t[\eps(\ee\times\bb)_i]$. 剩余 $\eps(\ee\times\partial_t\bb)_i$ 由法拉第定律化成 $\eps[(\ee\cdot\nabla)E_i-\frac12\partial_iE^2]$. 所以
\[
f_i+\partial_tg_i=\partial_jT_{ij},\qquad
\vct g=\eps\ee\times\bb=\vct S/c^2,
\]
\[
T_{ij}=\eps(E_iE_j-\tfrac12\delta_{ij}E^2)
+\frac1{\mu_0}(B_iB_j-\tfrac12\delta_{ij}B^2).
\]
对固定体积 $\Omega$ 积分, 首先得到
\[
\int_\Omega\vct f\dd^3x+\frac{\dd}{\dd t}\int_\Omega\vct g\dd^3x
=\oint_{\partial\Omega}\vct T\cdot\nn\dd a.
\]
若边界包围同一物质系统, 无物质穿越边界且没有其他外力, 第一项才等于 $\dd\vct P_\mathrm{mech}/\dd t$, 于是总动量的变化率等于应力边界积分. 一般固定控制体还须计入机械动量的流入流出和其他机械应力. 这里 $T_{ij}$ 表示第 $j$ 方向边界上的第 $i$ 分量力. 按本符号约定, 向外传出的场动量通量张量为 $-T_{ij}$, 不能把应力与外向动量流混为相同符号.

对平面波沿 $+z$ 传播, $T_{zz}=-u$, 确认它的外向动量流为 $u$. 吸收面的辐射压力为 $\langle S\rangle/c$, 完全反射为 $2\langle S\rangle/c$.

\originalsection{角动量}
场角动量密度为 $\rr\times\vct g$. 对动量守恒式乘 $\epsilon_{k\ell i}x_\ell$ 并积分, 分部积分的附加体项为 $\epsilon_{kji}T_{ij}=0$, 因为 $T$ 对称. 在上一节固定物质系统且无其他外力矩的条件下, 物质与场的总角动量只通过边界的应力矩交换; 有物质穿越固定边界时还须加入机械角动量通量. 静态 $\ee$ 与 $\bb$ 也可储存非零场动量, “静态”只表示场不随时间变化.

\begin{example}
半径 $R$ 的球壳均匀带电 $Q$, 以角速度 $\omega\hat{\vct z}$ 缓慢转动. 已知球内磁场为 $B_0\hat{\vct z}$, $B_0=\mu_0Q\omega/(6\pi R)$. 求场角动量.
\end{example}
\begin{solution}
球内电场为零, 只有球外贡献.

外场为 $E_r=Q/(4\pi\eps r^2)$, $B_\theta=B_0R^3\sin\theta/(2r^3)$. 所以 $g_\phi=QB_0R^3\sin\theta/(8\pi r^5)$, 轴向角动量密度为 $r\sin\theta g_\phi$. 用 $\int\sin^2\theta\dd\Omega=8\pi/3$ 和 $\int_R^\infty r^{-2}\dd r=1/R$, 得
$L_z=QB_0R^2/3=\mu_0Q^2R\omega/(18\pi)$.
\end{solution}
\originalsection{电阻为何从侧面接收场能}
设一段长直圆柱导线半径 $a$、长度 $\ell$, 导电率 $\sigma$, 直流电流 $I$ 沿 $+z$. 在远离端点的均匀电流近似中, $E_z=I/(\sigma\pi a^2)$; 导线侧面切向电场连续, 表面磁场 $B_\phi=\mu_0I/(2\pi a)$. 所以侧面的 $\vct S=\ee\times\bb/\mu_0$ 指向径向内侧, 因 $\hat{\vct z}\times\hat{\vct\phi}=-\hat{\vct r}$.

单位时间从侧面进入的总能量为
$2\pi a\ell\,E_zB_\phi/\mu_0=I E_z\ell=I^2\ell/(\sigma\pi a^2)$.
体内焦耳功率积分也给
$\int\jj\cdot\ee\dd^3x=\sigma E_z^2\pi a^2\ell=I^2R_\mathrm{wire}$,
$R_\mathrm{wire}=\ell/(\sigma\pi a^2)$. 静态场并不意味着没有能流; 这里 $\partial_tu=0$, 入流立即转成热. 维持电流的外电路负责建立上述场, 端部及返回导线附近需用完整电路几何处理. 这个局部计算解释本科作业9只给教材题号的相关概念, 不冒充缺失的教材题面.


本章对应 8.07 作业10、期末题1、2及 8.311 作业1、2, 保持真空场动量与介质动量问题的范围区分.
